\documentclass{article}
\begin{document}
\section*{Задача 1. Найти производную $f(x) = x^3 + 2x^2 - 5x + 7$}

Применяю правило дифференцирования суммы и степенную формулу
$\frac{d}{dx} x^n = n x^{n-1}$:

\begin{align*}
f'(x) &= \frac{d}{dx}\left(x^3\right) + \frac{d}{dx}\left(2x^2\right) - \frac{d}{dx}\left(5x\right) + \frac{d}{dx}\left(7\right) \\
      &= 3x^2 + 4x - 5 + 0 \\
      &= 3x^2 + 4x - 5
\end{align*}

\textbf{Ответ:} $f'(x) = 3x^2 + 4x - 5$.

\section*{Задача 2. Производная $g(x) = \sin(x) \cdot \cos(x)$}

Использую правило произведения:
\begin{align*}
g'(x) &= \cos(x) \cdot \cos(x) + \sin(x) \cdot (-\sin(x)) \\
      &= \cos^2(x) - \sin^2(x) \\
      &= \cos(2x)
\end{align*}

\textbf{Ответ:} $g'(x) = \cos(2x)$.
\end{document}
